% ===== sec/01_introduction.tex =====
\section{Introduction}
\label{sec:introduction}

\IEEEPARstart{T}{he} contemporary literature on autonomous software agents is rich in architectural diagrams but short on formal definitions.  An agent is variously described as ``a system that perceives and acts'', ``a planner that calls tools'', or ``an LLM in a loop'', and each definition is stated at a level of abstraction that admits neither a decision procedure nor a complexity bound.  The consequence is that two agents with radically different computational power and reasoning capabilities are routinely compared as if they belonged to a single class, and verification claims about agentic systems are correspondingly difficult to falsify.

We are not the first to ask whether the Chomsky hierarchy is the right scaffold for thinking about this problem.  In particular, Koohestani et al.\ pose the question \emph{Are agents just automata?}~\cite{koohestani2025automata} and answer affirmatively, establishing a memory-based equivalence between simple-reflex agents and finite automata, hierarchical task-decomposition agents and pushdown automata, and unrestricted-memory agents and thus Turing machines~\cite{turing1936computable}.  

Although the answer provided by Koohestani et al.\ is generally acceptable, the question they pose is the preliminary one that requires more profound and detailed thinking.  The question that does practical work is operational: \emph{given that agents are automata, how can the Chomsky type of one agent be used to \textbf{bound} the Chomsky type of another?}  In fact, this is known as the producer--consumer problem of Dijkstra~\cite{dijkstra1968cooperating} re-stated in the language of formal grammars.  The present paper offers a formal methodology that answers the posed question in the affirmative, with formal verification techniques inferred into the consumer side and the idealised theoretical apparatus of automata built into the producer side.

\subsection{Why a narrow definition}

We deliberately keep the definition of an agent narrow, taking
only what is forced upon us by the theory of automata Turing
introduced in 1936~\cite{turing1936computable} and the language
hierarchy Chomsky introduced in
1956--59~\cite{chomsky1956three,chomsky1959certain}.  We define an agent as {\it ``an effectful machine that, when started on an
admissible input, emits a finite string of operation symbols and
then halts (or fails to halt)''}.  The collection of all such
strings, taken across the input domain, is the agent's
\emph{trace language}.  Two agents with the same trace language
are, by stipulation, the same agent.  Whatever cannot be read off
from the trace language is, for the purposes of this paper, not
treated as an agent property.

The motivation for keeping the definition narrow is twofold.
First, it gives a direct line from the agent to the Chomsky
hierarchy and to the recognising automata.  An agent whose trace
language is regular (Type-3) is realisable by a finite-state
controller and admits decidable verification in linear
time~\cite{rabin1959finite}.  An agent whose trace language is
context-free (Type-2) is realisable by a pushdown controller and
admits $O(n^3)$ membership~\cite{hopcroft2006automata}.  An agent
whose trace language is recursively enumerable (Type-0) is, in
the worst case, \emph{undecidable to verify} --- a fact whose
modern statement we inherit from Robinson's resolution of
Hilbert's tenth
problem~\cite{robinson1949definability,davis1961decision}, which
established that recursively enumerable sets are exactly the
Diophantine sets and that the membership problem for them admits
no general algorithm.

Second, the narrow definition gives an empirical handle.  If the
agent is its trace language, then the trace language can be
recorded, stored, and compared; experiments produce strings, and
those strings can be classified by the smallest grammar class
that generates them.  This places the verification and validation of agentic systems on
the same footing as the verification and validation of compilers, parsers, and protocol
implementations~\cite{aho1972theory,hoare1969axiomatic} ---
domains in which the formal-language framing has, for half a
century, been the foundation of decidable correctness.

%\subsection{Finite Automata Formulation of Agents's Behavior: The Goal-Driven Producer and Data-Driven Consumer}
\subsection{The Goal-Driven Producer and Data-Driven Consumer}

Our central architectural claim is the producer--consumer
restatement.  There is exactly one \emph{Goal-Driven} producer
agent.  Its trace language is Type-0: recursively enumerable
and, in general, undecidable.  In practice it is realised as a
large language model wrapped in a tool-calling loop, and the
loop is permitted to enumerate any finite-length sequence of
tool invocations from the union of the sub-agents' alphabets.
The Goal-Driven producer has no a priori knowledge of \emph{what}
the goal is; it knows only the specification of the input domain
and the tool catalogue.

There is exactly one \emph{Data-Driven} consumer agent, the
\emph{verifier}.  Its trace language is Type-3: regular and, by
Rabin and Scott's construction~\cite{rabin1959finite}, decidable
in linear time.  The verifier is concretely realised as the
deterministic finite automaton of an Aho--Corasick keyword
index~\cite{aho1975corasick} compiled from the competition's
empirical fixtures: the unit-test corpus, the submission-validity
tests, the column-name checks, and the integrity tests. % on the \texttt{train.csv}\,/\,\texttt{test.csv}\,/\,\texttt{sample\_submission.csv} triple.
The binding is the entirety of our architectural claim:
\begin{quote}
\itshape A single Goal-Driven Type-0 producer, tamed by a single
Data-Driven Type-3 consumer, results in success.
\end{quote}


Type-2 and Type-1 sub-agents are not designed; they are
``\emph{derived}''.  They appear as the canonical projections of the
recursively enumerated trace onto the parts of the orchestration
substrate that admit, respectively, a stack discipline (the
back-edges of the planner/developer/reviewer loop, the
recursive \texttt{retry\_phase\_on\_low\_score}) and a
linear-bounded workspace (the deterministic Aho--Corasick scanner
and the LBA-style \texttt{compute\_dependency\_graph} that walks
the phase-to-phase dependency relation).  We make these
projections precise in Section~\ref{sec:chomsky}.

A subtlety: nothing in the framework prevents a Type-3 agent from
itself acting as a producer, in which case it backward-infers
constrained changes onto the tape (e.g., updates to the keyword
set $K$ as new task fixtures arrive).  Such a Type-3 producer is
admissible \emph{provided} the inference terminates in a
decidable verdict.  When it does not, the classical halting
problem of Turing~\cite{turing1936computable} re-asserts itself
on the Type-3 producer's loop and the architecture loses its
verification budget.  We discuss the engineering consequence in
Section~\ref{sec:discussion}: in the realised system, decidability
is enforced by the regular envelope of the verifier and, where
the data-driven side begins to drift, by a final human
verification and validation step that closes the loop.

\subsection{The Start Language and its Recursive Enumeration}

The system is initialised with a single \emph{start language}
$L_0$. %, given concretely in \texttt{data/trace\_language\_original.csv}.  
The start language
contains twelve sub-agents and their seed action alphabets ---
nothing else.  The Goal-Driven Type-0 producer, bound by the
Data-Driven Type-3 consumer, then \emph{recursively enumerates}
$L_0$ onto the \textsc{Spaceship Titanic} task instance.  The
result is the trace language $L_1$. %, recorded in\texttt{data/trace\_language\_kaggle.csv}.  
$L_1$ contains
\emph{thirteen} sub-agents.  The thirteenth,
\texttt{kaggle\_trainer}, is the emergent agent: it does not
exist in $L_0$ and is materialised \emph{by the act of
enumeration}, after the verifier first refuses traces that omit
the training step.

We claim that this enumeration is the operational meaning of
``Type-0 recursive enumeration'' in the agentic setting.
Turing's machines enumerate strings; Robinson's theorem connects
the enumerated sets to Diophantine arithmetic; our agent
enumerates tool-call strings, and the connection between the
enumerated strings and the empirical Kaggle competition is
mediated by the Aho--Corasick verifier.

\subsection{Idealisation and physical bounding}

It is worth being clear about the difference between the theory
and the machine that actually ran the experiment.  The
Goal-Driven agent's trace language is \emph{idealised} Type-0.
The actual execution was bounded by the finite tape and finite
RAM of a single laptop and, for the heavier ablation runs, by
the High Performance Computing Center at Texas Tech University.
The realised machine is, formally, a (very large) linear-bounded
automaton --- Type-1.  The idealisation is sound for our
purposes: every trace producible by the realised LBA is also
producible by the idealised Turing machine, and the verified
trace language is, by closure under intersection at the
verifier's Chomsky level (Theorem~\ref{thm:closure}), bounded
above by Type-3.  The \emph{physical} computation is therefore
far more restricted than the theoretical envelope, and any
verification claim we discharge about the idealised system holds
a fortiori on the bounded machine.  This is the standard
idealisation of
Hartmanis and Stearns~\cite{hartmanis1965computational}, and we
adopt it on those terms.

\subsection{Contributions}

This paper makes the following key contributions:

\begin{enumerate}
\item A formal definition of the trace language $L(\mathcal{A})$
of an agent $\mathcal{A}$ as the language generated by an
effectful controller over a finite operation alphabet, together
with a proof that the resulting object satisfies the Chomsky
hierarchy verbatim (Section~\ref{sec:trace}).

\item A producer--consumer formulation to the question posed by 
Koohestani et al.\:  a single Goal-Driven
Type-0 producer, intersected with a single Data-Driven Type-3
consumer realised as an Aho--Corasick DFA, suffices to produce a
verified trace language whose Chomsky class is bounded above by
the consumer's class (Section~\ref{sec:chomsky},
Theorem~\ref{thm:closure}).
\item A demonstration that Type-2 and Type-1 sub-agents need not
be designed but can be \emph{derived} as projections of the
recursively enumerated trace onto stack-discipline and
bounded-workspace substrates respectively
(Section~\ref{sec:methodology}).
\item A discussion of how a Type-3 agent may itself act as a
producer, backward-inferring constrained changes onto the tape,
under the necessary condition that decidability be preserved
(Section~\ref{sec:discussion}); failing this condition, the
halting problem returns and the verification budget collapses.
\item An empirical instantiation on the \textsc{Spaceship Titanic}
task from \textsc{MLE-bench} achieving placement $710/2330$
($\approx\!30.5\%$, top tertile).  The start language
\texttt{trace\_language\_original.csv}, the recursively
enumerated trace \texttt{trace\_language\_kaggle.csv}, and the
notebooks that produced the submission are released at\\
\url{https://www.kaggle.com/sweeden-ttu/competitions/spaceship-titanic}\\
(Sections~\ref{sec:expsetup}--\ref{sec:results}).
\end{enumerate}

\subsection{Roadmap}

Section~\ref{sec:related} surveys related work, opening with
Koohestani et al.\ and the producer--consumer framing.
Section~\ref{sec:preliminaries} fixes notation drawn from
Turing's automata theory and Chomsky's hierarchy.
Section~\ref{sec:trace} gives the rigorous definition of the
trace language.  Section~\ref{sec:chomsky} classifies trace
languages by Chomsky type and proves the closure theorem.
Section~\ref{sec:indexed} situates Aho's indexed grammars and
the Aho--Corasick verifier in the framework.
Section~\ref{sec:methodology} describes the Goal-Driven\,/\,
Data-Driven binding and the recursive enumeration of the start
language.  Sections~\ref{sec:expsetup}--\ref{sec:results} report
the experimental study and result on \textsc{Spaceship Titanic}.
Section~\ref{sec:discussion} discusses idealisation, physical
bounding, and the Type-3-as-producer subtlety.
Section~\ref{sec:conclusion} concludes.
